\documentclass[11pt]{article}
\usepackage{preamble}
\setlength{\parskip}{4pt}

\begin{document}

\daybanner{AP Statistics \textperiodcentered\ Unit 4 \textperiodcentered\ Topic 4.2 \textperiodcentered\ B: 2027-01-14 \textperiodcentered\ E: 2027-03-01 \textperiodcentered\ TEACHER KEY}{Two Interfaces, Twelve Specialists}

\sectionbanner{CED TETHER}

{\small
\begin{itemize}[leftmargin=1.5em,itemsep=2pt,topsep=2pt]
  \item \textbf{Skill 3.C} --- Describe probability distributions.
  \item \textbf{Skill 3.D} --- Construct a confidence interval, provided conditions for inference are met.
  \item \textbf{Skill 1.D} --- Identify an appropriate inference method for confidence intervals.
  \item \textbf{Skill 4.C} --- Verify that inference procedures apply in a given situation.
  \item \textbf{EU VAR-7} --- The t-distribution may be used to model variation.
  \item \textbf{LO VAR-7.A} --- Describe t-distributions. [Skill 3.C]
  \item \textbf{EK VAR-7.A.1} --- When s is used instead of $\sigma$ to calculate a test statistic, the corresponding distribution, known as the t-distribution, varies from the normal distribution in shape, in that more of the area is allocated to the tails of the density curve than in a normal distribution.
  \item \textbf{EK VAR-7.A.2} --- As the degrees of freedom increase, the area in the tails of a t-distribution decreases.
  \item \textbf{EU UNC-4} --- An interval of values should be used to estimate parameters, in order to account for uncertainty.
  \item \textbf{LO UNC-4.O} --- Identify an appropriate confidence interval procedure for a population mean, including the mean difference between values in matched pairs. [Skill 1.D]
  \item \textbf{EK UNC-4.O.1} --- Because $\sigma$ is typically not known for distributions of quantitative variables, the appropriate confidence interval procedure for estimating the population mean of one quantitative variable for one sample is a one-sample t-interval for a mean.
  \item \textbf{EK UNC-4.O.2} --- For one quantitative variable, X, that is normally distributed, the distribution of t = ($\bar{x}$ - $\mu$)/(s/$\surd$n) is a t-distribution with n - 1 degrees of freedom.
  \item \textbf{EK UNC-4.O.3} --- Matched pairs can be thought of as one sample of pairs. Once differences between pairs of values are found, inference for confidence intervals proceeds as for a population mean.
  \item \textbf{LO UNC-4.P} --- Verify the conditions for calculating confidence intervals for a population mean, including the mean difference between values in matched pairs. [Skill 4.C]
  \item \textbf{EK UNC-4.P.1} --- In order to calculate confidence intervals to estimate a population mean, we must check for independence and that the sampling distribution is approximately normal: a. To check for independence: (i) Data should be collected using a random sample or a randomized experiment. (ii) When sampling without replacement, check that n $\leq$ 10\% N.
  \item b. To check that the sampling distribution of $\bar{x}$ is approximately normal (shape): (i) If the observed distribution is skewed, n should be greater than 30. (ii) If the sample size is less than 30, the distribution of the sample data should be free from strong skewness and outliers.
  \item \textbf{LO UNC-4.Q} --- Determine the margin of error for a given sample size for a one-sample t-interval. [Skill 3.D]
  \item \textbf{EK UNC-4.Q.1} --- The critical value t* with n - 1 degrees of freedom can be found using a table or computer-generated output.
  \item \textbf{EK UNC-4.Q.2} --- The standard error for a sample mean is given by SE = s/$\surd$n, where s is the sample standard deviation.
  \item \textbf{EK UNC-4.Q.3} --- For a one-sample t-interval for a mean, the margin of error is the critical value (t*) times the standard error (SE), which equals t*(s/$\surd$n).
  \item \textbf{LO UNC-4.R} --- Calculate an appropriate confidence interval for a population mean, including the mean difference between values in matched pairs. [Skill 3.D]
  \item \textbf{EK UNC-4.R.1} --- The point estimate for a population mean is the sample mean, $\bar{x}$.
  \item \textbf{EK UNC-4.R.2} --- For the population mean for one sample with unknown population standard deviation, the confidence interval is $\bar{x}$ $\pm$ t*(s/$\surd$n).
\end{itemize}
}
\textbf{Essential question:} How do linked measurements choose an interval and bound a causal claim?\par
\textbf{DOK-3 skill:} 4.C \textperiodcentered\ \textbf{FRQ pattern:} {\small\texttt{paired-\allowbreak{}vs-\allowbreak{}independent-\allowbreak{}interval-\allowbreak{}and-\allowbreak{}scope}}\par
\textbf{DOK rationale:} Part (c) selects a procedure from the pairing, uses interval and SRS scope, diagnoses fixed-order confounding, and traces the independence and causal consequences.\par

\frameworkphaseheader{Do Now (first take) $\rightarrow$ Explore (video + follow-along) $\rightarrow$ Exit (finish + turn in)}{1 $\rightarrow$ 3}{45}{%
\begin{itemize}[leftmargin=1.5em,itemsep=2pt,topsep=2pt]
  \item Board slide up at the bell. Ask whether two columns of measurements necessarily mean two independent samples.
  \item Run the follow-along from minute 5 to 33; return to require procedure and scope statements.
  \item Collect at the door. Score part (c) only, E/P/I.
\end{itemize}
}{%
\begin{itemize}[leftmargin=1.5em,itemsep=2pt,topsep=2pt]
  \item Choose a procedure, complete the follow-along, finish (a)--(c), revise, and turn in.
\end{itemize}
}{%
\begin{itemize}[leftmargin=1.5em,itemsep=2pt,topsep=2pt]
  \item What is one observational unit here: a time or a specialist?
  \item What changed along with the interface between measurements?
\end{itemize}
}{Solo teacher. Tie the procedure to person-level differences and separate change from causation.}

\begin{callout}{\IconWarn\ WATCH FOR}{calloutred}
\begin{itemize}[leftmargin=1.5em,itemsep=2pt,topsep=2pt]
  \item Counting 24 independent observations or using 23 degrees of freedom.
  \item Checking separate-interface distributions instead of the 12 differences.
  \item Treating an interval above zero as proof that no practice or time effect exists.
\end{itemize}
\end{callout}

\sectionbanner{SCORING PART (c) --- E / P / I}

\textbf{Expected elements}
\begin{itemize}[leftmargin=1.5em,itemsep=2pt,topsep=2pt]
  \item \textbf{selects-paired-procedure} (required) --- Selects Priya and treats 12 linked differences as one sample
  \item \textbf{uses-interval-evidence} (required) --- Uses the positive interval $(0.9507,3.0493)$ under the protocol
  \item \textbf{bounds-population-scope} (required) --- Uses the SRS and 10 percent check to limit population scope
  \item \textbf{bounds-causation} (required) --- Uses fixed order and no comparison group to bound causation
  \item \textbf{states-two-reader-consequences} (required) --- Diagnoses false independence and false causal isolation
\end{itemize}

{\small\renewcommand{\arraystretch}{1.25}
\noindent\begin{tabular}{|p{0.5in}|p{5.9in}|}
\hline
\textbf{E} & Chooses matched pairs, uses the interval, bounds both scopes, and diagnoses both errors. \\ \hline
\textbf{P} & Chooses matched pairs but incompletely explains interval scope, confounding, or lost pairing. \\ \hline
\textbf{I} & Treats the columns as independent, counts 24 units, or claims the interval proves causation. \\ \hline
\end{tabular}}

\clearpage
\focusbanner{TODAY'S PROBLEM --- annotated}

Suppose a hospital randomly samples 12 of its 180 billing specialists. Each completes the same task with the current interface and, one week later, the redesign; everyone uses the current interface first and there is no comparison group. Let $d=\text{current}-\text{redesigned}$ time in seconds. The differences are shown.\par\smallskip \textbf{Priya:} ``Use a matched-pairs $t$ interval for $\mu_d$; fixed order limits a causal claim.''\par \textbf{Mateo:} ``Two columns require an independent two-sample interval; a positive interval proves the redesign caused faster work.''\par
\begin{center}
\begin{tikzpicture}
\begin{axis}[
  width=0.68\linewidth,
  height=1.28in,
  ymin=0, ymax=5, ytick=\empty, axis y line=none, axis x line=bottom,
  xlabel={Current time minus redesigned time (seconds)}, tick label style={font=\small}, label style={font=\small}
]
\addplot[only marks, mark=*, mark size=1.9pt, color=dayblue] coordinates {
  (-1,1)
  (0,1)
  (1,1)
  (1,2)
  (2,1)
  (2,2)
  (2,3)
  (2,4)
  (3,1)
  (3,2)
  (4,1)
  (5,1)
};
\end{axis}
\end{tikzpicture}
\end{center}
\begin{firsttakebox}{FIRST TAKE (5 min)}
Whose procedure is more defensible, Priya's or Mateo's? Cite how the two measurements are linked.\par
\answer{Not graded. Mateo's two-column rule is intentionally plausible; students should revise it by identifying the observational unit and the rival explanation created by fixed order.}
\end{firsttakebox}

\textit{Rules callout ``MATCH THE INTERVAL TO THE PAIRS (keep on the page)'' is printed on the student sheet.}\par\medskip

\dokbadge{1}~\textbf{(a)} Give the number of observational units and pairs, define $\mu_d$, state $df$, and interpret a positive difference.\par
\answer{There are 12 specialists and 12 linked pairs, not 24 independent units. Here $\mu_d$ is the population mean current-minus-redesigned time for the hospital's specialists under this fixed sequence; $df=11$. A positive $d$ means faster work at the later redesigned measurement.}

\dokbadge{2}~\textbf{(b)} Check random sampling, 10 percent, and the shape of the differences. Compute $\bar d$, $s_d$, SE, and a 95 percent interval for $\mu_d$; interpret it.\par
\answer{The SRS satisfies $12\leq0.10(180)=18$. The difference dotplot has no strong skew or outliers. Since the values sum to 24 and squared deviations from 2 sum to 30, $\bar d=2.000$, $s_d=\sqrt{30/11}=1.6514$, and $SE=1.6514/\sqrt{12}=0.4767$. With $t^*_{11}=2.201$, the interval is $2.000\pm1.0493=(0.9507,3.0493)$. We are 95 percent confident the population mean current-minus-redesigned time for the 180 specialists under this sequence is between 0.9507 and 3.0493 seconds.}

\dokbadge{3}~\textbf{(c)} Choose the warranted analysis from the linked data. Use the interval for the 180 specialists, then use fixed order and no comparison group to bound causation. Explain the two consequences of the rejected procedure and causal wording.\par
\answer{Priya is warranted because the 12 within-specialist differences form one sample. An independent analysis discards that linkage and uses the wrong standard error. The positive interval $(0.9507,3.0493)$ supports a positive mean change under this sequence, and the SRS supports generalization to the 180 specialists. But current always came first and there was no comparison group, so practice or time is confounded with interface. The interval does not isolate causation; Mateo's account falsely suggests 24 independent people and a causal redesign effect.}

\begin{summaryexitbox}{EXIT}
One thing the video changed about my first take: \hrulefill
\end{summaryexitbox}

\end{document}
